\section*{Chapter \ref{ch:b2:acyl-substitution} --- Carboxylic Acid Derivatives}

\begin{solution}{exo:b2:acyl-substitution:1}
Ethanamide $<$ ethyl ethanoate $<$ ethanoic anhydride $<$ ethanoyl chloride.
\end{solution}

\begin{solution}{exo:b2:acyl-substitution:2}
Ethanoic acid (+ \ce{HCl}); ethyl ethanoate; ethanamide (+ ammonium chloride, the second equivalent
of ammonia trapping \ce{HCl}); ethanoic anhydride (+ \ce{NaCl});
\textit{N},\textit{N}-dimethylethanamide (+ triethylammonium chloride).
\end{solution}

\begin{solution}{exo:b2:acyl-substitution:3}
\ce{C6H5COOCH3 + NaOH -> C6H5COONa + CH3OH}. $13.6/136.15 = \qty{0.0999}{mol}$, as much
\ce{NaOH}: \qty{4.00}{g}.
\end{solution}

\begin{solution}{exo:b2:acyl-substitution:4}
$\gamma$-Butyrolactone (oxolan-2-one): a five-membered ring of four carbons and one oxygen.
\end{solution}

\begin{solution}{exo:b2:acyl-substitution:5}
Ammonia adds to the carbonyl carbon; a proton moves from the nitrogen to the oxygen (or to the
leaving group); ethoxide (as ethanol, after proton transfer) leaves and the \ce{C=O} reforms:
ethanamide and ethanol.
\end{solution}

\begin{solution}{exo:b2:acyl-substitution:6}
The phenolic \ce{OH} is acylated (acid catalysis): 2-acetoxybenzoic acid; by-product ethanoic acid.
\end{solution}

\begin{solution}{exo:b2:acyl-substitution:7}
2-Methylpropan-2-ol. The first addition gives propanone, more electrophilic than the ester, which
takes the second methyl at once.
\end{solution}

\begin{solution}{exo:b2:acyl-substitution:8}
Butane-1,4-diol: the ester is reduced to two primary alcohols, the ring being opened.
\end{solution}

\begin{solution}{exo:b2:acyl-substitution:9}
One equivalent: $x^2/(1 - x)^2 = 4$, $x = 2/3$. Three equivalents: $x^2/[(1 - x)(3 - x)] = 4$, $3x^2 - 16x
+ 12 = 0$, $x = 0.90$. Removing the water as it forms (Dean--Stark) drives it further.
\end{solution}

\begin{solution}{exo:b2:acyl-substitution:10}
\ce{SOCl2} gives benzoyl chloride; with two equivalents of diethylamine (or one plus triethylamine)
it gives the amide. Acid and amine mixed directly undergo an acid--base reaction: diethylammonium
benzoate, which gives the amide only on strong heating.
\end{solution}

\begin{solution}{exo:b2:acyl-substitution:11}
Three KOH per triolein: $3 \times 56.11/885.4 = \qty{0.190}{g}$ per gram, a \omterm{def:b2:acyl-substitution:saponification}{saponification} value of
190.
\end{solution}

\begin{solution}{exo:b2:acyl-substitution:12}
\ce{NaBH4}: only the ketone, to a secondary alcohol. \ce{LiAlH4}: the ketone to the secondary
alcohol, the ester to a primary alcohol (and releases the alcohol part), the amide to an amine.
\end{solution}

\begin{solution}{pb:b2:acyl-substitution:1}
\textbf{1.} Oleic acid \ce{C18H34O2}; glycerol \ce{C3H8O3}.
\textbf{2.} \ce{C57H104O6}: $57 \times 12.011 + 104 \times 1.008 + 6 \times 15.999 = \qty{885.45}{g/mol}$.
\textbf{3.} Three.
\textbf{4.} The cis double bond kinks the chains, which pack poorly: weak intermolecular forces, a
low melting point.
\textbf{5.} Glycerol and three sodium oleate, a soap.
\textbf{6.} A mixture of methyl esters of fatty acids.
\textbf{7.} \ce{C57H104O6 + 3CH3OH -> C3H8O3 + 3C19H36O2}.
\textbf{8.} The methoxide ion, formed from methanol and the base, is a far stronger nucleophile than
methanol; it attacks the ester carbonyls.
\textbf{9.} The carbonyl carbon of one ester group bearing \ce{O-}, the \ce{OCH3} of methoxide and the
glyceryl oxygen: tetrahedral.
\textbf{10.} \omterm{def:b2:acyl-substitution:transesterification}{Transesterification} is an equilibrium with a constant near 1; excess methanol shifts it
towards the methyl esters.
\textbf{11.} Glycerol separates as a lower layer: removing a product drives the equilibrium.
\textbf{12.} No: methoxide is regenerated in each exchange.
\textbf{13.} It neutralises the base, giving sodium oleate (soap): the catalyst is consumed.
\textbf{14.} It saponifies them: hydroxide (from water and methoxide) cuts the esters to soap and
methanol, irreversibly.
\textbf{15.} It consumes base, emulsifies the mixture and prevents the glycerol layer from separating.
\textbf{16.} By an acid-catalysed \omterm{def:b2:acyl-substitution:esterification}{esterification} with methanol first, which turns the free acids into
methyl esters.
\textbf{17.} Water hydrolyses esters and converts the catalyst into hydroxide, which saponifies.
\textbf{18.} $10^6/885.45 = \qty{1129}{mol}$.
\textbf{19.} $3 \times 1129 \times 32.04 = \qty{108.6}{kg}$.
\textbf{20.} $3 \times 1129 \times 296.50 = \qty{1004.6}{kg}$.
\textbf{21.} $1000 + 108.6 = \qty{1108.6}{kg}$ in; $104.0 + 1004.6 = \qty{1108.6}{kg}$ out.
\textbf{22.} \qty{20}{kg} of oleic acid, $20\,000/282.47 = \qty{70.8}{mol}$, as much \ce{NaOH}:
\qty{2.83}{kg}.
\textbf{23.} $1129 \times 92.09 = \boldsymbol{\approx \qty{104}{kg}}$ of glycerol per tonne.
\end{solution}
